\section{Applications to image processing and statistics}
A data matrix has one observation in each \emph{column} and one measured variable in each row. Its mean column is $M=N^{-1}\sum_{k=1}^N X_k$. Subtract $M$ from every column to form $B$; then the sample covariance matrix is $S=BB^T/(N-1)$, for $N>1$.

For a unit coefficient vector $c$, the scalar index $y=c^TX$ has variance $c^TSc$. The greatest possible variance is the largest eigenvalue of $S$, attained at a corresponding unit eigenvector. Principal component scores normally use the centered observation $X-M$. Centering changes an index by a constant and hence does not change its variance. Signs of principal component directions are arbitrary: reversing a unit eigenvector reverses every score but preserves its explained variance.

\practice{7.5.P1}{Center the weight-height data and calculate covariance}
The five-column mean is $M=(130,65)^T$. Subtracting it from every observation gives
\[B=\begin{bmatrix}-10&-5&-5&5&15\\-4&-5&-1&3&7\end{bmatrix}.\]
Each row sums to zero, checking the centering step. The cross-product is
\[BB^T=\begin{bmatrix}400&190\\190&100\end{bmatrix}.\]
Since the sample has $N=5$, divide by $N-1=4$:
\[S=\frac14BB^T=\begin{bmatrix}100&\frac{95}{2}\\\frac{95}{2}&25\end{bmatrix}.\]
Its diagonal entries are the weight and height sample variances; its off-diagonal entry is their sample covariance.

\practice{7.5.P2}{A one-number size index}
Diagonalizing the covariance matrix gives variances approximately $123.01859218$ and $1.98140782$. A unit direction for the larger one is
\[c\approx\begin{bmatrix}0.899901\\0.436094\end{bmatrix}.\]
Thus a centered size score is $y\approx0.899901(w-130)+0.436094(h-65)$. Its variance is the larger eigenvalue. The total variance is $\operatorname{tr}S=125$, so the fraction retained is $0.98414874$, or about $98.4149\%$. This index uses the original units specified by the question; standardizing the two variables to unit variance would define a different PCA problem.

\exercise{7.5.1}{Center all observations and form the sample covariance}
\step{Compute the mean of each row using all six observations.}
\[M=\begin{bmatrix}12\\10\end{bmatrix}.\]
\step{Subtract that column from every observation.}
\[B=\begin{bmatrix}7&10&-6&-9&-10&8\\2&-4&-1&5&3&-5\end{bmatrix}.\]
Each row sums to zero. This checks that no observation was omitted and that the rows and columns were not reversed.
\step{Form the sample covariance with denominator $6-1$.}
\[BB^T=\begin{bmatrix}430&-135\\-135&80\end{bmatrix},\qquad S=\frac15BB^T=\begin{bmatrix}86&-27\\-27&16\end{bmatrix}.\]
For example, the off-diagonal entry is the sum of six products of paired centered measurements, divided by $5$. The two diagonal entries are sums of squares divided by $5$, so they are nonnegative as variances must be.

\exercise{7.5.2}{Center all observations and form the sample covariance}
\step{Compute the mean of each row using all six observations.}
\[M=\begin{bmatrix}4\\9\end{bmatrix}.\]
\step{Subtract that column from every observation.}
\[B=\begin{bmatrix}-3&1&-2&2&3&-1\\-6&2&-3&-1&6&2\end{bmatrix}.\]
Each row sums to zero. This checks that no observation was omitted and that the rows and columns were not reversed.
\step{Form the sample covariance with denominator $6-1$.}
\[BB^T=\begin{bmatrix}28&40\\40&90\end{bmatrix},\qquad S=\frac15BB^T=\begin{bmatrix}\frac{28}{5}&8\\8&18\end{bmatrix}.\]
For example, the off-diagonal entry is the sum of six products of paired centered measurements, divided by $5$. The two diagonal entries are sums of squares divided by $5$, so they are nonnegative as variances must be.

\exercise{7.5.3}{Find both principal component directions}
Use the exact covariance matrix from Exercise 1:
\[S=\begin{bmatrix}86&-27\\-27&16\end{bmatrix}.\]
Its characteristic equation is $t^{2} - 102 t + 647=0$. In decreasing order its eigenvalues are approximately
\[\lambda_1=95.20407221,\qquad\lambda_2=6.79592779.\]
Choose the following orthonormal eigenvector columns, in that same order:
\[P\approx\begin{bmatrix}0.946515&0.322659\\-0.322659&0.946515\end{bmatrix}.\]
The two principal component scores are the entries of $Y=P^T(X-M)$, with mean
\[M=\begin{bmatrix}12\\10\end{bmatrix}.\]
Their covariance matrix is diagonal, $P^TSP=\operatorname{diag}(\lambda_1,\lambda_2)$, so they are uncorrelated. The first score explains about $93.33733\%$ of total variance; the second explains the remainder. A sign-reversed column gives an equivalent principal component.

\exercise{7.5.4}{Find both principal component directions}
Use the exact covariance matrix from Exercise 2:
\[S=\begin{bmatrix}\frac{28}{5}&8\\8&18\end{bmatrix}.\]
Its characteristic equation is $\frac{5 t^{2} - 118 t + 184}{5}=0$. In decreasing order its eigenvalues are approximately
\[\lambda_1=21.92126474,\qquad\lambda_2=1.67873526.\]
Choose the following orthonormal eigenvector columns, in that same order:
\[P\approx\begin{bmatrix}0.440130&0.897934\\0.897934&-0.440130\end{bmatrix}.\]
The two principal component scores are the entries of $Y=P^T(X-M)$, with mean
\[M=\begin{bmatrix}4\\9\end{bmatrix}.\]
Their covariance matrix is diagonal, $P^TSP=\operatorname{diag}(\lambda_1,\lambda_2)$, so they are uncorrelated. The first score explains about $92.88672\%$ of total variance; the second explains the remainder. A sign-reversed column gives an equivalent principal component.

\exercise{7.5.5}{A leading principal component from a supplied covariance matrix}
The full covariance matrix needed for this calculation is supplied in the exercise:
\[S=\begin{bmatrix}\frac{4103}{25}&\frac{3273}{100}&\frac{2026}{25}\\\frac{3273}{100}&\frac{13486}{25}&\frac{24913}{100}\\\frac{2026}{25}&\frac{24913}{100}&\frac{18911}{100}\end{bmatrix}.\]
No original satellite image is required to diagonalize this matrix. Compute its symmetric eigenproblem and normalize the eigenvector of largest eigenvalue. The eigenvalues, in decreasing order, are approximately
\[(677.49712437,\ 181.58112062,\ 33.59175501).\]
A maximizing unit direction is
\[c\approx\begin{bmatrix}0.129554\\0.874423\\0.467547\end{bmatrix}.\]
Thus $y_1=c_1\widehat x_1+c_2\widehat x_2+c_3\widehat x_3$, where the hats denote mean-deviation variables, is the first principal component. The uncentered combination has the same variance. The maximum variance is approximately $677.49712437$, and total variance is $\operatorname{tr}S=\frac{89267}{100}$. Therefore the explained fraction is
\[\frac{\lambda_1}{\operatorname{tr}S}\approx 0.75895586,\quad\text{or }75.89559\%.\]
The actual computation used the real symmetric eigenvalue routine:
\begin{verbatim}
eigenvalues,P = np.linalg.eigh(S)
c = P[:,-1]
print(c,eigenvalues[-1],
      100*eigenvalues[-1]/np.trace(S))
\end{verbatim}
The maximum eigenvalue is the last one because \texttt{eigh} returns increasing order. A reversed eigenvector sign is equivalent. Orthogonality and the residual $Sc-\lambda_1c$ were checked at full precision before rounding.

\exercise{7.5.6}{A leading principal component from a supplied covariance matrix}
The full covariance matrix needed for this calculation is supplied in the exercise:
\[S=\begin{bmatrix}\frac{741}{25}&\frac{919}{50}&5\\\frac{919}{50}&\frac{1041}{50}&\frac{703}{50}\\5&\frac{703}{50}&\frac{2921}{100}\end{bmatrix}.\]
No original satellite image is required to diagonalize this matrix. Compute its symmetric eigenproblem and normalize the eigenvector of largest eigenvalue. The eigenvalues, in decreasing order, are approximately
\[(51.69566979,\ 24.53510921,\ 3.43922100).\]
A maximizing unit direction is
\[c\approx\begin{bmatrix}0.615525\\0.599424\\0.511683\end{bmatrix}.\]
Thus $y_1=c_1\widehat x_1+c_2\widehat x_2+c_3\widehat x_3$, where the hats denote mean-deviation variables, is the first principal component. The uncentered combination has the same variance. The maximum variance is approximately $51.69566979$, and total variance is $\operatorname{tr}S=\frac{7967}{100}$. Therefore the explained fraction is
\[\frac{\lambda_1}{\operatorname{tr}S}\approx 0.64887247,\quad\text{or }64.88725\%.\]
The actual computation used the real symmetric eigenvalue routine:
\begin{verbatim}
eigenvalues,P = np.linalg.eigh(S)
c = P[:,-1]
print(c,eigenvalues[-1],
      100*eigenvalues[-1]/np.trace(S))
\end{verbatim}
The maximum eigenvalue is the last one because \texttt{eigh} returns increasing order. A reversed eigenvector sign is equivalent. Orthogonality and the residual $Sc-\lambda_1c$ were checked at full precision before rounding.

\exercise{7.5.7}{A maximal-variance scalar combination}
For the data in Exercise 1, the covariance matrix is
\[S=\begin{bmatrix}86&-27\\-27&16\end{bmatrix}.\]
Choose the unit coefficient direction
\[c\approx\begin{bmatrix}0.946515\\-0.322659\end{bmatrix}.\]
The requested uncentered variable is $y_1\approx (0.946515)x_1+(-0.322659)x_2$. Its variance is $c^TSc=\lambda_1\approx 95.20407221$, the maximum by the quadratic-form theorem. Using $x_1-M_1,x_2-M_2$ instead would simply subtract its mean and leave that variance unchanged.
The total variance is $\operatorname{tr}S=102$, so the explained percentage is approximately $93.33733\%$. Both the unit-length constraint and the covariance eigenvector equation were checked before rounding the coefficients.

\exercise{7.5.8}{A maximal-variance scalar combination}
For the data in Exercise 2, the covariance matrix is
\[S=\begin{bmatrix}\frac{28}{5}&8\\8&18\end{bmatrix}.\]
Choose the unit coefficient direction
\[c\approx\begin{bmatrix}0.440130\\0.897934\end{bmatrix}.\]
The requested uncentered variable is $y_1\approx (0.440130)x_1+(0.897934)x_2$. Its variance is $c^TSc=\lambda_1\approx 21.92126474$, the maximum by the quadratic-form theorem. Using $x_1-M_1,x_2-M_2$ instead would simply subtract its mean and leave that variance unchanged.
The total variance is $\operatorname{tr}S=\frac{118}{5}$, so the explained percentage is approximately $92.88672\%$. Both the unit-length constraint and the covariance eigenvector equation were checked before rounding the coefficients.

\exercise{7.5.9}{A unit eigenvector maximizes the index variance}
The sample variance of $y=c^TX$ is $c^TSc$, with $c^Tc=1$. The characteristic roots of the supplied symmetric matrix are $9,\ 6,\ 3$, with multiplicity. Its largest eigenvalue is $9$. Solving $(S-(9)I)c=0$ and normalizing gives
\[c=\begin{bmatrix}\frac{1}{3}\\\frac{2}{3}\\\frac{2}{3}\end{bmatrix}.\]
The exact checks are $c^Tc=1$ and $Sc=(9)c$. Hence the index
\[y=(\frac{1}{3})x_1 + (\frac{2}{3})x_2 + (\frac{2}{3})x_3\]
has the largest possible variance, $9$. This accounts for the fraction $\frac{1}{2}$ of the total variance $18$.

\exercise{7.5.10}{A unit eigenvector maximizes the index variance}
The sample variance of $y=c^TX$ is $c^TSc$, with $c^Tc=1$. The characteristic roots of the supplied symmetric matrix are $15,\ 3,\ 3$, with multiplicity. Its largest eigenvalue is $15$. Solving $(S-(15)I)c=0$ and normalizing gives
\[c=\begin{bmatrix}\frac{\sqrt{6}}{6}\\\frac{\sqrt{6}}{3}\\\frac{\sqrt{6}}{6}\end{bmatrix}.\]
The exact checks are $c^Tc=1$ and $Sc=(15)c$. Hence the index
\[y=(\frac{\sqrt{6}}{6})x_1 + (\frac{\sqrt{6}}{3})x_2 + (\frac{\sqrt{6}}{6})x_3\]
has the largest possible variance, $15$. This accounts for the fraction $\frac{5}{7}$ of the total variance $21$.

\exercise{7.5.11}{Covariance identities and their dimensions}
\step{(a) The transformed mean remains zero.}
Write $B=[X_1\ \cdots\ X_N]$ and $C=[Y_1\ \cdots\ Y_N]=P^TB$. Since the original mean is zero, $\sum_kX_k=0$. Thus
\[
\frac1N\sum_kY_k=P^T\left(\frac1N\sum_kX_k\right)=0.
\]
Equivalently, $B\mathbf1=0$ implies $C\mathbf1=P^TB\mathbf1=0$. No orthogonality of $P$ is required.
\step{(b) Compute the new covariance using part (a).}
Because both matrices are already centered,
\[
S_Y=\frac1{N-1}CC^T
=\frac1{N-1}P^TBB^TP=P^TSP.
\]
The transpose on $C^T$ reverses factors and gives $C^T=B^TP$, which explains the rightmost $P$.

\exercise{7.5.12}{Covariance identities and their dimensions}
For orthogonal $P$, the relation $X=PY$ is equivalent to $Y=P^TX$. Centering commutes with that linear change of variables, and Exercise 11 gives covariance $S_Y=P^TSP$. Total variance is the trace. The cyclic trace identity yields
\[
\operatorname{tr}(S_Y)=\operatorname{tr}(P^TSP)
=\operatorname{tr}(SPP^T)=\operatorname{tr}(SI)=\operatorname{tr}(S).
\]
Thus orthogonal changes redistribute variance among new coordinate directions without creating or destroying its total. This proof also covers data not originally centered, because covariance first subtracts the transformed mean.

\exercise{7.5.13}{Covariance identities and their dimensions}
Let $M=N^{-1}\sum_{k=1}^NX_k$ and $B=[X_1-M\ \cdots\ X_N-M]$. Multiply the column partition of $B$ by the row partition of $B^T$:
\[
BB^T=\sum_{k=1}^N(X_k-M)(X_k-M)^T.
\]
Hence the sample covariance is
\[
\boxed{S=\frac1{N-1}\sum_{k=1}^N(X_k-M)(X_k-M)^T.}
\]
Every summand is a $p\times p$ outer-product matrix. Its $(i,j)$ entry is the product of the centered measurements of variables $i$ and $j$ in observation $k$. When $p=1$, the outer product is just the square $(t_k-m)^2$, giving exactly the usual scalar sample-variance formula. The sum has $N$ terms, one per observation; the lowercase $n$ printed as an upper bound in the displayed scalar formula should be read as that sample size.
